\section{Where the time goes}
\label{sec:analysis}

\generated{table_main_breakdown}

\generated{table_ablation}

\generated{table_ablation_speedup}

The breakdown separates the cost of one pivot from the number of pivots. The
first is bounded by the structures: full pricing scans the $n-1$ non-basic
variables of the residual, the direction visits the entering subtree and the
positive component, and the ratio test examines every basic slack of the
residual, so a pivot costs $O(n+m)$ regardless of how small the change to the
forest is. The second is a property of the instance and of the entering rule,
and it is the quantity the ablation is designed to expose.

The ablation switches off one implementation choice at a time with respect to the
frozen configuration, on the validation manifest, with the same data and the same
acceptance thresholds as the main campaign.

\subsection{What separates the easy families from the hard ones}

It is tempting to attribute the cost to the number of blocks, since each one
costs a fresh oracle with a fresh basis. The measurements do not support that.
Across generated instances of two thousand vertices, all with the frozen
configuration: a chain produces $11$ blocks and $4\,640$ pivots in $0.116$\,s;
a sparse DAG of out-degree three produces $282$ blocks and $134\,830$ pivots in
$1.311$\,s; a layered DAG of twenty levels produces $232$ blocks and
$348\,205$ pivots in $6.079$\,s; and a layered DAG of eight levels with dense
connections between them produces only $85$ blocks --- fewer than either sparse
instance --- yet $528\,841$ pivots, and does not finish within the limit. More
blocks is not more work.

The quantity that separates them is pivots \emph{per oracle}, expressed as a
multiple of $n$: about $0.21n$ on the chain, $0.24n$ on the sparse DAG,
$0.75n$ on the twenty-level instance and $3.1n$ on the eight-level dense one.
The layered structure gives every entering variable many precedence slacks
between vertices that are all at zero, so each oracle spends thousands of
zero-step pivots walking to its optimum.

Two consequences follow for where effort should go. Reusing the basis across
residual calls, which would save one initialization per oracle, cannot help the
hard case: a hundred initializations against six hundred thousand pivots. What
can help is either making each pivot cheaper --- incremental pricing and updates
confined to the modified paths, since pricing and rebuilding are together four
fifths of the solver's time on this family --- or attacking the number of
degenerate pivots directly with a lexicographic or perturbed ratio test. Only the
second changes the multiplier, and it is the one with no implementation here.

This also explains the shape of the comparison with the parametric backend. That
backend pays nothing for degeneracy: it solves one maximum flow per subproblem on
disjoint vertex sets. Its advantage is largest exactly on the family where the
forest's pivot count per oracle is largest, which is not a coincidence of
constants but a difference in what each method spends its work on.

\subsection{The dual simplex}
\label{sec:dual-results}

\generated{table_dual}
\generated{table_dual_speedup}

The measurements above say that the primal loop is dominated by zero-step
pivots. Primal degeneracy and dual degeneracy are different defects, and here
they are wildly asymmetric: the primal kind is a per-pivot phenomenon, at a
median of $99.3\%$ of pivots over the whole dual campaign, while the only source
of the dual kind is a tie between block ratios, which is a per-block event and
therefore bounded by the number of blocks. That asymmetry is the standard
indication for the dual method, and it was implemented and measured.

Two structural facts make it unusually clean on this problem. The vertex variables
are never primal infeasible, since their value is $1/w(P)$ on the positive
component and zero elsewhere, so every primal infeasibility is a precedence
slack, that is an arc leaving $P$; and the basis is primal feasible exactly when
$P$ is closed, which is therefore the stopping rule. The starting basis --- every
vertex its own tree, $P$ the vertex of maximum ratio --- is dual feasible by
construction and costs $O(n)$.



The prediction is confirmed on the quantity it was made about, and refuted on
every other. Table~\ref{tab:dual} and table~\ref{tab:dual-speedup} report the
dedicated campaign on the frozen validation split, $31$ instances in two
replicas under the same protocol as every other campaign.

Degeneracy collapses: the primal sits at a median of $99.3\%$ of pivots with
zero step, the dual at $0.0\%$. Nothing else improves. The dual performs
\emph{more} pivots, not fewer --- a median of $7\,813$ against the primal's
$3\,253$ on the instances both solve --- and it is slower: the paired speedup
against the primal has a median of $0.90$, ranging from $0.34$ to $2.47$. It
also solves less: $56$ of $62$ runs against $62$ of $62$, losing all three
replicas of a dense DAG of two thousand vertices that the primal solves in
about eleven seconds. The variant that picks the largest tree is worse again,
median $0.78$ and $54$ of $62$.

So the reduction of degeneracy was real and worth nothing. The dual traded null
steps for a longer path, which is the outcome that the primal's own diagnosis
should have predicted: the cost was never the null steps as such, it was the
number of pivots needed to walk to the optimum, and the dual walks farther. On
the same instances the parametric backend is $3.89$ times the primal at the
median and up to $135$ times on one of them.

Choosing the leaving arc by a structural criterion instead of by smallest index
does not help either, and there is a structural reason: every violated slack has
the same value $-1/w(P)$, so a most-infeasible rule has nothing to discriminate.
The largest-tree rule is the one carried into the campaign, and it is the slowest
of the three settings measured.

Correctness is not at stake in any of this. The dual returns the same canonical
blocks as the primal on the whole exhaustive suite and on fifteen hundred
random instances checked with the diagnostics enabled at every pivot, and it
produces its certificate through the same verified code. If dual feasibility is
lost to rounding, the solver restarts from the primal initialization rather than
inheriting an infeasible basis, so the answer is certified either way.

